\section{Introduction}

The automated discovery of physical laws from data is a central challenge in AI for Science~\citep{wang2023scientific,raghu2020survey}. While traditional Symbolic Regression (SR) methods like Genetic Programming (e.g., PySR) are effective, they struggle with combinatorial search spaces and lack semantic priors \cite{cranmer2020discovering}. Large Language Models (LLMs) offer a promising alternative by leveraging pre-trained knowledge of physics and code generation~\citep{chen2021evaluating,wei2022chain}. However, applying LLMs to equation discovery faces a critical barrier: \textit{instruction brittleness}~\citep{sclar2024quantifying,mizrahi2024state}---the phenomenon where small changes in prompt phrasing lead to dramatically different outputs.

Standard frameworks such as LLM-SR \citep{shojaee2024llmsr} and LLM4ED \citep{du2024llm4ed} typically employ fixed prompt templates (e.g., ``Find the equation that fits this data''). This rigid approach mimics a ``static researcher'' unable to adapt their line of questioning. An LLM prompted with ``Find the simplest equation'' might over-regularize and miss important terms; the same model prompted with ``Find the most accurate equation'' might hallucinate spurious terms or fixate on memorized but irrelevant formulas. When faced with complex nonlinear dynamics, static prompts often lead models into local optima from which they cannot escape. This paper makes a key observation: \textbf{the optimal instruction depends on the current state of the generation process}. Early in the search, exploratory prompts (``propose novel functional forms'') are beneficial; later, refinement prompts (``prune redundant terms'') become essential. A fixed prompt cannot capture this state-dependency. 

To overcome this, we introduce \ours, a framework that treats the prompt instruction not as a fixed input, but as a \textit{dynamic hyperparameter} to be optimized during the search. We construct a discrete space of reasoning strategies and employ Bayesian Optimization (BO) to navigate this space, enabling the system to actively switch strategies based on the error profile of current candidates. Our contributions are: \textbf{(1) Dynamic Instruction Tuning}: We formalize prompt engineering in scientific discovery as a discrete Bayesian Optimization problem, enabling adaptive control over the LLM's generation mode. \textbf{(2) Sample Efficiency}: Unlike Reinforcement Learning approaches to prompt optimization, our BO-based method is highly sample-efficient, making it feasible for computationally expensive scientific evaluations. \textbf{(3)Empirical Robustness}: We demonstrate that \textsc{NeuroSym-BO} achieves higher recovery rates on benchmark PDEs, solving cases where static prompting fails.
% \begin{itemize}
%     \item \textbf{Dynamic Instruction Tuning}: We formalize prompt engineering in scientific discovery as a discrete Bayesian Optimization problem, enabling adaptive control over the LLM's generation mode.
%     \item \textbf{Sample Efficiency}: Unlike Reinforcement Learning approaches to prompt optimization, our BO-based method is highly sample-efficient, making it feasible for computationally expensive scientific evaluations.
%     \item \textbf{Empirical Robustness}: We demonstrate that \textsc{NeuroSym-BO} achieves higher recovery rates on benchmark PDEs, solving cases where static prompting fails.
% \end{itemize}